\plannedchapter{perf、PMU 和 Top-Down 分析}
{本章讲硬件性能计数器如何把“猜测瓶颈”变成证据。}
{\item cycles、instructions、IPC 能说明什么，不能说明什么？
\item Top-Down 的 Retiring、Bad Speculation、Frontend Bound、Backend Bound 如何解读？
\item WSL2 为什么可能无法使用完整 PMU？}
{\item PMU event、sampling、multiplexing。
\item \code{perf stat}、\code{perf record/report}。
\item Top-Down methodology。}
{在裸机或可用环境下对前面 kernel 做 perf stat/report；在 WSL2 记录限制并用替代证据分析。}
